\section{Tổng quan về đề tài}

\subsection{Lý do chọn đề tài}

Trong bối cảnh giáo dục hiện đại, việc rèn luyện tư duy logic và khả năng giải quyết vấn đề cho học sinh tiểu học ngày càng được chú trọng. Toán tư duy không chỉ giúp các em nâng cao năng lực học thuật mà còn là nền tảng quan trọng cho các kỹ năng như phân tích, suy luận và sáng tạo. Tuy nhiên, các hình thức học truyền thống thường khô khan, thiếu tính trực quan và dễ làm giảm hứng thú của học sinh.

Bên cạnh đó, công nghệ giáo dục (EdTech) đang phát triển mạnh mẽ, mở ra cơ hội ứng dụng các phương pháp mới giúp học sinh tiếp cận kiến thức một cách sinh động và hiệu quả hơn. Việc xây dựng một trang web học toán tư duy với giao diện trực quan, hoạt động tương tác và mini-game sinh động sẽ giúp các em vừa học vừa chơi, từ đó tăng khả năng ghi nhớ và phát triển tư duy tự nhiên.

Chính vì vậy, nhóm quyết định lựa chọn đề tài “Website học toán tư duy cho học sinh tiểu học” nhằm tạo ra một công cụ học tập hiện đại, dễ sử dụng, phù hợp với lứa tuổi và có khả năng hỗ trợ việc học tập một cách hiệu quả và thú vị hơn. Đề tài hướng tới việc thu hẹp khoảng cách giữa lý thuyết toán học và ứng dụng thực tiễn thông qua các trải nghiệm số hóa.

\subsection{Thách thức hiện tại}

Việc triển khai các nền tảng giáo dục toán tư duy tại Việt Nam hiện nay vẫn đang đối mặt với nhiều khó khăn, xuất phát từ cả phương pháp tiếp cận lẫn hạ tầng công nghệ. Qua quá trình khảo sát thực trạng học tập của học sinh tiểu học, nhóm đã tổng hợp được các thách thức cốt lõi như sau:

\begin{itemize} \item \textbf{Thiếu môi trường học tập trực quan và tương tác:} Phần lớn tài liệu toán tư duy hiện nay ở dạng sách in hoặc bài tập truyền thống, khiến học sinh khó hình dung, đặc biệt với các bài đòi hỏi suy luận hoặc thao tác trực tiếp. \item \textbf{Khó duy trì hứng thú học tập:} Nội dung học không được game hóa hoặc thiết kế hấp dẫn khiến học sinh dễ chán nản, không duy trì được sự tập trung trong thời gian dài. \item \textbf{Thiếu hệ thống theo dõi tiến trình học tập:} Giáo viên và phụ huynh khó nắm bắt được mức độ hiểu bài của học sinh, tiến trình học hoặc điểm mạnh – điểm yếu để có hướng hỗ trợ kịp thời. \item \textbf{Chưa có công cụ học toán tư duy phù hợp độ tuổi:} Nhiều nền tảng hiện nay thiết kế cho lứa tuổi lớn hơn, sử dụng thuật ngữ phức tạp hoặc giao diện khó sử dụng đối với học sinh tiểu học. \item \textbf{Hạn chế trong việc cá nhân hóa trải nghiệm học:} Mỗi học sinh có khả năng tiếp thu khác nhau, nhưng các tài liệu hiện tại chưa đáp ứng tốt việc tùy chỉnh độ khó hoặc cung cấp gợi ý phù hợp năng lực từng em. \end{itemize}

Những rào cản này không chỉ ảnh hưởng đến hiệu quả tiếp thu của học sinh mà còn tạo ra áp lực cho phụ huynh trong việc kèm cặp. Do đó, việc xác định rõ các thách thức này là cơ sở quan trọng để nhóm đề xuất các tính năng giải quyết triệt để vấn đề trong phần thiết kế hệ thống.

\subsection{Mục tiêu của đề tài}

Mục tiêu tổng quát của đề tài là phát triển một hệ thống học toán tư duy trực tuyến dành riêng cho học sinh tiểu học, giúp các em rèn luyện khả năng suy luận và tư duy logic thông qua các bài học tương tác. Hệ thống được kỳ vọng sẽ tạo ra một môi trường học tập cá nhân hóa, nơi năng lực của từng học sinh được ghi nhận và phát triển theo lộ trình phù hợp. Để đạt được điều đó, đề tài tập trung vào các mục tiêu cụ thể sau:

\begin{itemize} \item Xây dựng các bài học toán tư duy với nội dung trực quan, dễ hiểu, bám sát chương trình giáo dục phổ thông mới của Bộ Giáo dục và Đào tạo Việt Nam. \item Phát triển bộ mini–game tương tác đa dạng trên nền tảng Cocos Creator 3.8, giúp học sinh phát triển tư duy thông qua các thao tác kéo thả, giải đố và tìm quy luật. \item Thiết lập hệ thống Gamification toàn diện bao gồm điểm thưởng, huy hiệu (badges) và bảng xếp hạng để thúc đẩy động lực tự học dựa trên các khung lý thuyết về trò chơi hóa. \item Xây dựng công cụ quản lý dành cho phụ huynh, cho phép theo dõi chi tiết tiến trình học tập, biểu đồ tăng trưởng kỹ năng và thời gian tương tác của con trẻ trên hệ thống. \end{itemize}

Việc hoàn thành các mục tiêu này sẽ mang lại một giải pháp giáo dục toàn diện, không chỉ dừng lại ở việc cung cấp kiến thức mà còn hình thành thói quen tư duy độc lập cho học sinh. Đây là tiền đề để xây dựng một cộng đồng học tập trực tuyến lành mạnh và bền vững.

\subsection{Phạm vi của đề tài}

Để đảm bảo tính khả thi và chất lượng hoàn thiện của hệ thống, đề tài tập trung nghiên cứu và triển khai trong những giới hạn cụ thể. Các nội dung thực hiện được phân định rõ ràng qua bốn khía cạnh chính bao gồm chức năng, kỹ thuật, người dùng và các phần nằm ngoài phạm vi nghiên cứu:

\begin{itemize} \item \textbf{Phạm vi chức năng:} Hệ thống tập trung cung cấp kho bài học tương tác, các mini-game rèn luyện phản xạ toán học và hệ thống báo cáo kết quả tự động. Các tính năng về quản lý tài khoản phụ huynh và cơ chế vinh danh thông qua huy hiệu cũng là trọng tâm của giai đoạn phát triển này.
	
	\item \textbf{Phạm vi kỹ thuật:} Đề tài ứng dụng các công nghệ hiện đại như Next.js cho giao diện người dùng, Node.js và MongoDB cho quản trị dữ liệu backend. Đặc biệt, phần bài tập tương tác được xây dựng chuyên sâu bằng engine Cocos Creator 3.8 để tối ưu hóa trải nghiệm trên trình duyệt.
	
	\item \textbf{Phạm vi người dùng:} Đối tượng phục vụ chính là học sinh tiểu học (từ lớp 1 đến lớp 5) và các phụ huynh tại Việt Nam có nhu cầu sử dụng công cụ học tập bổ trợ bằng tiếng Việt.
	
	\item \textbf{Phạm vi ngoài đề tài:} Trong giai đoạn này, hệ thống sẽ không hỗ trợ tính năng thi đấu trực tuyến thời gian thực (Real-time Battle) giữa nhiều người chơi, không phát triển ứng dụng di động riêng biệt (Native App) và không cung cấp nội dung cho cấp học trung học.
\end{itemize}